% secctions/3_visual.tex
\section{Composición del Modo Matemático Avanzado e Inserción de Recursos Gráficos}

\subsection{Fundamentos del Entorno Matemático en \LaTeX}
La maquetación tipográfica de expresiones matemáticas constituye una de las mayores fortalezas estructurales de \LaTeX, diferenciándolo cualitativamente de cualquier procesador de textos convencional. El sistema procesa las fórmulas dentro de un espacio abstracto gobernado por reglas métricas e isotrópicas distintas a las del texto regular, donde los espacios en blanco son ignorados y las variables se renderizan automáticamente en una tipografía cursiva matemática optimizada. Para activar este motor tipográfico, se dispone de dos modalidades operacionales:

\begin{itemize}
    \item \textbf{Modo Matemático en Línea (\textit{Inline Math Mode}):} Se utiliza para integrar expresiones matemáticas breves o variables individuales de forma fluida dentro del flujo continuo de un párrafo. Sintácticamente, se delimita encerrando la expresión entre símbolos de dólar simples (\texttt{\$... \$}). Por ejemplo, la ecuación de continuidad lineal $f(x) = x + y$ se acopla perfectamente a la línea sin alterar el interlineado vertical del párrafo.
    
    \item \textbf{Modo Matemático Desplegado (\textit{Display Math Mode}):} Se reserva para ecuaciones complejas, teoremas de alta jerarquía o demostraciones formalizadas que requieren un aislamiento visual. Este modo suspende el párrafo, centra la expresión en una línea independiente y optimiza las proporciones tipográficas de los operadores integrales o fraccionarios. Se invoca de manera estándar mediante el delimitador de doble símbolo de dólar (\texttt{\$\$... \$\$}). 
\end{itemize}

Para propósitos de referenciación cruzada formal en tesis o artículos científicos, se prefiere la instanciación del entorno cerrado \texttt{\textbackslash begin\{equation\}} y \texttt{\textbackslash end\{equation\}}, el cual activa el modo desplegado y, simultáneamente, asigna un identificador numérico secuencial automático al margen derecho del documento.

\subsection{Sintaxis de Estructuras Matemáticas Complejas}
La construcción de fórmulas en \LaTeX\ opera de forma anidada y declarativa, permitiendo representar nociones del cálculo infinitesimal y el álgebra con absoluta precisión visual.

\subsubsection{Subíndices, Superíndices y Fracciones}
Los componentes atómicos del álgebra se gestionan mediante operadores de posición. Los subíndices se declaran con el carácter de guion bajo (\texttt{\_}) y los superíndices o potencias con el acento circunflejo (\texttt{\^}). Cuando el argumento posicional excede un único carácter, la sintaxis exige el encapsulamiento estricto entre llaves argumentales; por ejemplo, la expresión $x_{i+1}^{2n}$ garantiza la correcta indexación del bloque. 

Por su parte, las fracciones se maquetan formalmente mediante la macro \\ \texttt{\textbackslash frac\{numerador\}\{denominador\}}, la cual ajusta dinámicamente el tamaño de la fuente interna para preservar la armonía de la línea en caso de anidamientos sucesivos.

\subsubsection{Operadores Acumulativos: Sumatorios, Productorios e Integrales}
La representación de límites, sumas infinitas e integrales se realiza mediante macros especializadas que admiten límites inferiores y superiores utilizando la misma sintaxis de subíndices y superíndices. En el modo desplegado, estos límites se posicionan matemáticamente arriba y abajo del operador principal de forma automática:
\begin{itemize}
    \item \textbf{Sumatorios y Productorios:} Invocados con \texttt{\textbackslash sum} y \texttt{\textbackslash prod} respectivamente, permiten estructurar series analíticas complejas como $\sum_{i=1}^{\infty} a_i$.
    \item \textbf{Cálculo Integral:} Definido mediante la macro \texttt{\textbackslash int}, permite estructurar integrales definidas de la forma $\int_{a}^{b} f(x)\,dx$. Es una buena práctica académica anteponer el comando de espaciado fino (\texttt{\textbackslash,}) antes del diferencial (\texttt{dx}) para asegurar una separación tipográfica óptima.
\end{itemize}

\subsubsection{Matrices y Álgebra Lineal}
A través del paquete core \texttt{amsmath}, \LaTeX\ prescinde de la complejidad de los entornos de alineación tabulares primitivos para el diseño de matrices. Se introducen entornos específicos que gestionan automáticamente los delimitadores perimetrales sin necesidad de declarar el número estricto de columnas:
\begin{itemize}
    \item \texttt{matrix}: Genera un arreglo bidimensional regular sin delimitadores visuales.
    \item \texttt{pmatrix}: Modela matrices estándar encapsuladas entre paréntesis formales.
    \item \texttt{bmatrix}: Diseña matrices indexadas estructuradas entre corchetes.
    \item \texttt{vmatrix}: Entorno estándar para la representación de determinantes mediante barras verticales puras.
\end{itemize}
La sintaxis interna respeta las convenciones del entorno de bloques, empleando el símbolo \texttt{\&} como separador de columnas y la doble barra invertida (\texttt{\textbackslash\textbackslash}) como operador de salto de fila.

\subsection{Inserción de Recursos Gráficos y Gestión del Entorno Flotante}
La incorporación de elementos visuales externos (como diagramas de arquitectura, capturas de datos o vectores en PDF) es regulada por el paquete \texttt{graphicx} en el preámbulo del sistema, el cual dota al documento de compatibilidad nativa con formatos estructurados como PNG, JPG y PDF vectorial.

\subsubsection{El Comando \textbackslash includegraphics y Parámetros Técnicos}
La importación física de un recurso se ejecuta mediante la instrucción \\ \texttt{\textbackslash includegraphics[opciones]\{ruta\}}. La macro provee un argumento opcional entre corchetes altamente versátil para realizar transformaciones geométricas directas sobre el lienzo, controlando variables críticas como:
\begin{itemize}
    \item \texttt{width}: Determina el ancho absoluto o relativo del gráfico (ej. \texttt{width=0.5\textbackslash textwidth} escala el elemento exactamente al 50\% del ancho disponible del texto bloque).
    \item \texttt{height}: Define la altura del recurso, permitiendo un escalado proporcional automático si se omite el ancho.
    \item \texttt{angle}: Aplica una rotación angular explícita en sentido antihorario expresada en grados enteros.
\end{itemize}

\subsubsection{Automatización de Leyendas, Etiquetas y Entornos Flotantes}
Debido a que los recursos gráficos y las tablas extensas constituyen bloques no segmentables, su posicionamiento arbitrario puede romper la continuidad estética de las páginas. Para resolver esto, \LaTeX\ introduce el concepto de \textbf{entornos flotantes} a través del bloque \texttt{figure}. Este entorno actúa como un contenedor dinámico que el motor de compilación reubica automáticamente en la sección superior (\texttt{t}), inferior (\texttt{b}) o en la posición exacta del código (\texttt{h}), optimizando el balance de espacios en blanco.

Dentro de este entorno, se estructuran tres comandos jerárquicos indispensables para la redacción universitaria:
\begin{enumerate}
    \item \texttt{\textbackslash caption\{...\}}: Define la leyenda explicativa o descripción técnica del recurso gráfico. \LaTeX\ indexa internamente esta macro para enumerar la figura correlativamente y habilitar su inclusión automática en el índice de ilustraciones (\texttt{\textbackslash listoffigures}).
    \item \texttt{\textbackslash label\{...\}}: Establece una meta-etiqueta alfanumérica única en la base de datos interna del proyecto. Es un requisito sintáctico estricto posicionar la macro \texttt{\textbackslash label} inmediatamente después de la declaración del \texttt{\textbackslash caption}.
    \item \texttt{\textbackslash ref\{...\}}: Invocado en el cuerpo del texto regular, permite automatizar las referencias cruzadas dinámicas (ej. \textit{"como se aprecia en la Figura \textbackslash ref\{etiqueta\}..."}), asegurando que si el orden de las imágenes cambia durante la edición, los números indexados en el texto se actualicen de manera exacta y transparente en la siguiente compilación.
\end{enumerate}